% ================================================================
% RESULTS SECTION
% ================================================================

This section presents the estimation results in four stages. Section~\ref{sec:gravity_results} establishes that the gravity framework fits bilateral film production flows. Section~\ref{sec:ppml_incentives} reports the PPML incentive estimates (Models~3--7). Section~\ref{sec:pair_fe_results} reports the fixed effects estimates (Models~8--10).. Section~\ref{sec:batis_results} compares these findings to the BaTIS robustness test results.

% ----------------------------------------------------------------
\subsection{The Gravity Framework for Film Production}
\label{sec:gravity_results}
% ----------------------------------------------------------------

Tables~\ref{tab:model1} and~\ref{tab:model2} report the baseline and full gravity specifications for both dependent variables. The gravity model fits bilateral film flows well, with pseudo-$R^2$ values of 0.62--0.66, though the determinants differ from goods trade in several notable respects. Origin GDP is consistently positive and significant, though modest for film counts (elasticity $\approx 0.05$) and much larger for budgets ($\approx 0.6$). Destination GDP is insignificant for film counts and significantly negative for budgets, inverting the standard gravity prediction. These results suggest that larger, wealthier origin countries fund more and higher-value overseas productions, but that their wealth means they have higher production costs such as wages which deter international productions.

Distance is negative and highly significant across all specifications, but the elasticity ($-0.24$ to $-0.28$ for film counts; $-0.18$ for budgets) is substantially smaller than the $-0.7$ to $-1.5$ range typically found in goods trade \citep{head2014}. This is consistent with the services trade literature, which finds that non-physical trade is less distance-sensitive \citep{walsh2006, kimura2006}. The smaller distance elasticity for budgets suggests that when distant productions do occur, they tend to be higher budget.

Common language is the single most important predictor of bilateral film flows, increasing production by approximately 69\% for film counts and 47\% for budgets (Model~2). This is larger than the average language effect of 44\% typically found in goods trade \citep{egger2012}, which makes sense given how language-intensive films are as a product. Contiguity increases film flows by 27\% for counts and 32\% for budgets, less than average good trade estimate of 50\% \citep{head2014}.

Colonial relationships are associated with approximately 26\% \textit{less} bilateral film production, contrary to the large positive colonial-tie effects found in goods trade. This is surprising but may reflect that post-colonial cultural differences mean that production in many ex-colonies has intentionally diverged from their colonial tradition (prominent examples include the Indian, Latin American, and Nigerian film industries). It may also reflect that film making is an industry which was relatively small in colonial times and has expanded significantly in recent decades, meaning it is less dependent on historical production networks. Regional trade agreements (RTA) are associated with 17--21\% less production, less than average good trade estimate of 50\% \citep{head2014}. A potential explanation for this is that RTA partners tend to have similar economic characteristics, so co-productions are unable to take advantage of differential factors such as lower wages. This is supported by the lower contiguity estimates compared to goods trade. Finally, economic freedom is insignificant across all specifications, which contrasts with \citet{kimura2006} where it was found to be a strong predictor of services trade.

% ----------------------------------------------------------------
\subsection{Incentive Effects: PPML Estimates}
\label{sec:ppml_incentives}
% ----------------------------------------------------------------

Tables~\ref{tab:model3}--\ref{tab:model7} report the PPML incentive estimates. Adding incentive variables does not destabilise the gravity controls, which remain essentially unchanged from Model~2.

\subsubsection{Aggregate Incentive Effects}

Model~3 (Table~\ref{tab:model3}) introduces the primary variables of interest. For film counts, the destination incentive coefficient implies that countries with incentive schemes attract 18\% more international productions, while the origin coefficient implies that studios reduce overseas filming by $12\%$ when their home country introduces incentives.

For budgets, an asymmetry emerges. The destination attraction effect is similar in magnitude ($+20\%$) but statistically insignificant, while the origin pull-back effect is much larger ($-38\%$) and statistically significant. Home-country incentives appear to most effectively retain high-budget productions, which have the most to gain in absolute terms from a domestic rebate on a large qualifying spend. This pattern of a strong destination effect for film counts but not budgets and a strong origin effect for budgets but not counts persists across Models~4--7.

\subsubsection{Incentives: Design, Generosity, Duration and Timing}

Model~4 (Table~\ref{tab:model4}) disaggregates destination incentives into three instrument types: cash rebates, tax credits, and tax shelters. The specification retains the aggregate destination incentive dummy, with cash rebates as the omitted reference category; the tax credit and tax shelter coefficients capture the differential effect of those instruments relative to a rebate scheme.

For film counts, cash rebates are the most effective instrument, and the estimate is similar to the Model~3 destination coefficient. Tax credits have a $+8.1\%$ differential and tax shelters have a negative $-9.3\%$, but these are both statistically insignificant. For budgets, cash rebates show a large positive coefficient ($+30\%$) but it is not statistically significant. Tax credits and tax shelts have negative differentials ($-8\%$ and $-67\%$ respectively) but both are insignificant. Tax shelters having a large negative differential is consistent with tax shelter regimes being structured around investor tax optimisation rather than genuine production attraction. However, given the small sample size (only 3 countries operate tax shelters regimes) as well as potential selection effects, we cannot draw strong conclusions from this result. This is true for tax credits too, which only cover 6 of the 57 countries with incentives and therefore doesn't constitute a robust sample.

Model~5 (Table~\ref{tab:model5}) replaces dummies with continuous generosity rates. Both the destination-side and origin-side effects scale roughly linearly with the rate, with each additional 10\% increase in the generosity rate for film counts and budgets roughly corresponding to a $+6--7\%$ destination effect and a $-7--15\%$ origin effect. At the mean active generosity of 27.6\%, the implied destination-effect is  $+17\%$ for film counts and $+21\%$ for budgets. The implied origin-effect is $-13\%$ for film counts and $+21\%$ for budgets. All estimates are statistically significant except for the budget destination-effect estimate, which is in line with the budget incentive effect also being insignificant.

Models~6 and~7 (Tables~\ref{tab:model6} and~\ref{tab:model7}) test whether incentive effects accumulate over time. Model~6 finds no evidence of duration effects: each additional year of incentive operation adds approximately $0.6\%$ to inward film production, but this is statistically insignificant. The benefit of having incentives appears to be a one-time level shift rather than a compounding advantage, providing no support to a ``cluster-building'' effect. However, Model~7 appears to show that \textit{when} a country adopted incentives matters. Early adopters (pre-2005) attract $26\%$ more films and $45\%$ more production spending, both highly significant. Late adopters attract only $15\%$ more films and show no statistically significant budget effect.  These results appear to show that incentives do not become more effective the longer they operate, but that early adopters' incentive schemes saw success because they entered a less competitive environment.

% ----------------------------------------------------------------
\subsection{Identification: Country-Pair Fixed Effects}
\label{sec:pair_fe_results}
% ----------------------------------------------------------------

Tables~\ref{tab:model8}--\ref{tab:model10} report the pair fixed effects (FE) OLS estimates, which absorb all time-invariant bilateral heterogeneity and identify effects exclusively from within-pair temporal variation.

\subsubsection{Pair FE Incentive Effects}

Model~8 (Table~\ref{tab:model8}) reports the core FE estimates. For film counts, the destination incentive effect survives at approximately $+9\%$, half of the PPML estimate of $+18\%$. This reduced is expected as pair FE strip out the cross-sectional correlation between incentive adoption and pre-existing production relationships. This quantifies the selection component, whereby countries with established film industries were more likely to adopt incentives, inflating the cross-sectional estimate. The portion that survives ($+9\%$) represents the genuine within-pair effect. This result is robust across subsamples. Excluding Anglo-to-Anglo pairs yields $+9.4\%$; excluding all Anglo countries yields $+7\%$; dropping EFW controls yields $+9.4\%$. The origin pull-back for film counts is $-9\%$ (compared to $-12\%$ estimate in Model~3) and significant in the full sample. The effect vanishes in the excluding-all-Anglo subsample, but the result is insignificant. Overall, both the destination and origin incentive effects are weaker in magnitude than under the PPML specifications but remain meaningful.

For budgets, the pair FE destination effect ($+22\%$) is large and significant in both the full sample and the Anglo-to-Anglo subsamples, although it disapears when all Anglo countries are excluded. This is in contrast to the Model~3 PPML estimate, which found a similar sized but statistically insignificant effect ($+20\%$). Meanwhile, the origin effect ($-13\%$) is insignificant across all pair FE subsamples, which is also in contrast to the large and statistically significant PPML estimate ($-38\%$). These results run counter to those found in Models~3--7, and overall strengthens the argument for a meaningful destination effect while undermining the case for an origin effect

\subsubsection{Event Study}

Model~9 (Table~\ref{tab:model9}) replaces the destination incentive dummy with relative-time dummies to examine dynamic effects and test parallel trends.

For film counts, the full-sample pre-treatment coefficients are small and statistically insigniciant, with only $\tau_{-5}$ marginally significant at the $10\%$ level. Post-treatment coefficients show a clear jump at $\tau_{+1}$ and $\tau_{+2}$, consistent with a one-to-two year lag between incentive introduction and production decisions responding, and estimates are consistently significant at the $5\%$ level for the first seven years. However, when all Anglo countries are excluded, the pre-treatment coefficients are larger and individually significant, with the post-treatment average (($10\%$)) falling \textit{below} the pre-treatment average (($15\%$)).

For budgets, the full-sample pre-treatment average is essentially zero ($1\%$), with no individually significant coefficients. Post-treatment estimates are small initially but several estimates from $\tau_{+5}$ onward are statistically significant at the $5\%$ and $10\%$ level.  The excluding-Anglo-to-Anglo subsample is the strongest specification: a slightly negative pre-treatment average ($-10\%$), followed by large and increasingly significant post-treatment effects. The subsample excludes all Anglo countries shows a similar result. Overall, the film count data shows strong evidence of a causal effect and this supported to a lesser extent by the budget data.

\subsubsection{Duration Under Pair FE}

Model~10 (Table~\ref{tab:model10}) confirms the null duration result from Model~6 under pair FE identification. The years-active coefficient is insignificant for film counts across all subsamples ($0.1\%--0.4\%$ per year) and for budgets in three of four subsamples. The one exception is the excluding-all-Anglo budget subsample, where years active is significantly \textit{negative} ($-6.9\%$ per year). This suggests that non-Anglo destinations may see incentive effectiveness decay as competitive entry erodes their initial advantage, but the small sample size (835 observations, 72 pairs) limits this result.

\subsubsection{Robust Estimator for Staggered Adoption}
\label{sec:results_bjs}

Table 12~\ref{tab:bjs} reports the FE estimates imputation estimator of \citet{borusyak2024}, which restricts the comparison group to never-treated and not-yet-treated units and is robust to the heterogeneous-treatment-effects bias that can affect two-way FE under staggered adoption.  Note that the imputation estimator does not naturally produce pre-treatment coefficients, since pre-treatment imputations are zero by construction. 

For film counts, the BJS results confirms Model~8 findings. The pooled Average Treatment Effect (ATT) in the full sample is $+11\%$, statistically significant, and close to the Model~8 estimate of $+9\%$.The dynamic pattern is cleaner than that in Model~9: the imputed effect at $\tau_{+0}$ is essentially zero, jumps to approximately $+8\%$ at $\tau_{+1}$, and grows to between $+11\%$ and $+17\%$ across years $\tau_{+2}$ through $\tau_{+8}$, with most post-treatment coefficients individually significant. Both Anglo subsamples yield a smaller pooled ATT of $+6\%$ but it remains statistically significant.

For budgets, the pooled ATT is $+9\%$ in the full sample but is insignificant and there is no coherent post-treatment trajectory. The no-Anglo pair and no-Anglo subsamples yield pooled ATT estimates of $-30\%$ and $-60\%$ respectively. These are notably highly negative and reserve the sign of the Model ~9 FE estimates, but only the no-Anglo estimate is statistically significant and it is dentified on a small sample (only 254 treated observations and 692 untreated observations). Nonetheless, these results suggest that weak positive budget effects in the earlier FE estimates were driven by heterogeneous-treatment-effects bias.

% ----------------------------------------------------------------
\subsection{Cross-Dataset Validation: BaTIS Audiovisual Services}
\label{sec:batis_results}
% ----------------------------------------------------------------

The parallel estimation on related audiovisual services data provides independent validation. The BaTIS results are reported in full in the appendix tables.

The gravity framework fits BaTIS more closely than IMDb in several respects. Both origin and destination GDP are large and highly significant, consistent with standard gravity predictions for services trade. Distance is larger and closer to goods trade estimates than the IMDb elasticities. Common language doubles audiovisual services flows ($+102\%$), exceeding even the largest IMDb effect ($+69\%$). Colonial ties and RTAs are both positive and significant in BaTIS, reversing the negative IMDb findings. This divergence likely reflects BaTIS capturing ongoing post-colonial media trade and the benefits of trade liberalisation for broader services flows, rather than discrete film production decisions.

The core incentive findings are similar. The destination incentive effect for BaTIS under pair FE is $+27\%$ (compared to $+45\%$ under BaTIS PPML), which is larger than the IMDb film count effect ($+9\%$) and also statistically significant. The BaTIS destination effect is stable across Anglo subsamples, providing stronger evidence than IMDb that the destination attraction effect generalises beyond the Anglo production network. The origin pull-back observed in IMDb PPML ($-12\%$) is also present in BaTIS PPML ($-15\%$), though under pair FE the BaTIS origin effect flips to \textit{positive} ($+22\%$) and is statistically significant. This suggests that the introduction of home-country incentives may simultaneously pull film shoots back while deepening broader audiovisual services relationships through expanded post-production and distribution activity.

BaTIS provides mixed evidence on duration effects. In PPML, years active is positive and significant ($+2.2\%$ per year) unlike the null IMDb result. However, under pair FE the coefficient reverses to significantly \textit{negative} ($-2.1\%$ per year), indicating that the PPML duration effect reflected cross-sectional selection rather than genuine within-pair accumulation.

The Model~7 event study is uninformative due to near-perfect multicollinearity among event-time indicators in the balanced panel structure \citep{borusyak2024}.